% ===============================================================
% ML-06d: Mi ciudad - Wiederholung (MUSTERLÖSUNG)
% Spanisch Klasse 7 - Sequenz 6: Mi ciudad
% ===============================================================

\documentclass[11pt, a4paper]{article}

\newif\ifswmodus
\swmodusfalse

\usepackage[utf8]{inputenc}
\usepackage[T1]{fontenc}
\usepackage[ngerman]{babel}
\usepackage{helvet}
\renewcommand{\familydefault}{\sfdefault}

\usepackage[margin=2cm]{geometry}
\usepackage{enumitem}
\usepackage{tabularx}
\usepackage{booktabs}

\usepackage{xcolor}
\usepackage{tcolorbox}
\tcbuselibrary{skins}

\usepackage{fancyhdr}
\usepackage{lastpage}
\usepackage{needspace}

% FARBEN
\definecolor{spanisch-color}{HTML}{c41e3a}
\definecolor{grau-color}{HTML}{888888}
\definecolor{hellgrau-color}{HTML}{f5f5f5}
\definecolor{orange-color}{HTML}{b35c00}
\definecolor{loesung-color}{HTML}{c62828}

\definecolor{sw-dark}{HTML}{000000}
\definecolor{sw-gray}{HTML}{666666}
\definecolor{sw-white}{HTML}{ffffff}

\ifswmodus
    \colorlet{spanisch}{sw-dark}
    \colorlet{grau}{sw-gray}
    \colorlet{hellgrau}{sw-white}
    \colorlet{orange}{sw-dark}
    \colorlet{loesung}{sw-dark}
\else
    \colorlet{spanisch}{spanisch-color}
    \colorlet{grau}{grau-color}
    \colorlet{hellgrau}{hellgrau-color}
    \colorlet{orange}{orange-color}
    \colorlet{loesung}{loesung-color}
\fi

\newcommand{\fachfarbe}{spanisch}

\newcommand{\thema}{Mi ciudad -- Wiederholung}
\newcommand{\sequenz}{06}
\newcommand{\block}{d}
\newcommand{\fach}{Spanisch}
\newcommand{\klasse}{7}
\newcommand{\maxpunkte}{20}
\newcommand{\datum}{\today}

\pagestyle{fancy}
\fancyhf{}
\renewcommand{\headrulewidth}{0.4pt}
\renewcommand{\footrulewidth}{0.4pt}

\lhead{\fach{} -- Klasse \klasse}
\chead{\textcolor{loesung}{\textbf{ML-\sequenz\block{} (MUSTERLÖSUNG)}}}
\rhead{\datum}

\lfoot{}
\cfoot{}
\rfoot{Seite \thepage\ von \pageref{LastPage}}

\newcommand{\punktebox}[1]{\hfill\fbox{\small \_/#1}}
\newcommand{\luecke}[1]{\rule{#1}{0.4pt}}
\newcommand{\lsg}[1]{\textcolor{loesung}{\textbf{#1}}}

\newtcolorbox{merkkasten}[1][]{
    colback=hellgrau,
    colframe=\fachfarbe,
    colbacktitle=white,
    coltitle=\fachfarbe,
    boxrule=1pt,
    arc=0pt,
    left=8pt, right=8pt, top=8pt, bottom=8pt,
    fonttitle=\bfseries,
    attach boxed title to top left={yshift=-2mm, xshift=5mm},
    boxed title style={boxrule=0pt, colback=white},
    title=#1
}

\newenvironment{aufgabe}[1]{%
    \needspace{5\baselineskip}%
    \nopagebreak[4]%
    \vspace{10pt}
    \noindent\textbf{\textcolor{\fachfarbe}{Aufgabe #1}}
    \nopagebreak[4]%
    \vspace{5pt}
    \nopagebreak[4]%
    \begin{list}{}{\leftmargin=0pt}
    \item[]
}{%
    \end{list}
}

\newcommand{\es}[1]{\textcolor{\fachfarbe}{\textbf{#1}}}

\begin{document}

\begin{center}
    {\Large\bfseries\textcolor{\fachfarbe}{Arbeitsblatt: \thema}}\\[0.3em]
    {\large\textcolor{loesung}{\textbf{--- MUSTERLÖSUNG ---}}}
\end{center}

\vspace{5pt}

\noindent
\begin{tabularx}{\textwidth}{|X|X|X|}
\hline
\textbf{Name:} \lsg{Musterschüler/in} &
\textbf{Klasse:} \klasse &
\textbf{Datum:} \datum \\
\hline
\end{tabularx}

\vspace{15pt}

% ===============================================================
% MERKKASTEN: Orte in der Stadt
% ===============================================================

\begin{merkkasten}[Merkkasten: Lugares en la ciudad (Orte in der Stadt)]

\begin{tabular}{ll|ll}
\es{la tienda} & = \lsg{das Geschäft} & \es{el supermercado} & = \lsg{der Supermarkt} \\[0.5em]
\es{el parque} & = \lsg{der Park} & \es{la farmacia} & = \lsg{die Apotheke} \\[0.5em]
\es{el cine} & = \lsg{das Kino} & \es{el hospital} & = \lsg{das Krankenhaus} \\[0.5em]
\es{la plaza} & = \lsg{der Platz} & \es{la estación} & = \lsg{der Bahnhof} \\[0.5em]
\end{tabular}

\end{merkkasten}

\vspace{10pt}

% ===============================================================
% MERKKASTEN: Verben ir und estar
% ===============================================================

\begin{merkkasten}[Merkkasten: Die Verben ir und estar]

\begin{tabular}{|l|l|l||l|l|l|}
\hline
\multicolumn{3}{|c||}{\textbf{ir} (gehen, fahren)} & \multicolumn{3}{c|}{\textbf{estar} (sein, sich befinden)} \\
\hline
yo & \lsg{voy} & voy & yo & \lsg{estoy} & estoy \\
\hline
tú & \lsg{vas} & vas & tú & \lsg{estás} & estás \\
\hline
él/ella & \lsg{va} & va & él/ella & \lsg{está} & está \\
\hline
nosotros & \lsg{vamos} & vamos & nosotros & \lsg{estamos} & estamos \\
\hline
\end{tabular}

\vspace{0.5em}
\textbf{Wichtig:} ir + a + el = ir \es{al} (z.B. Voy \es{al} cine.)

\end{merkkasten}

\vspace{10pt}

% ===============================================================
% MERKKASTEN: Wegbeschreibung
% ===============================================================

\begin{merkkasten}[Merkkasten: Wegbeschreibung]

\begin{tabular}{ll|ll}
\es{todo recto} & = \lsg{geradeaus} & \es{cerca de} & = \lsg{in der Nähe von} \\[0.5em]
\es{a la derecha} & = \lsg{nach rechts} & \es{lejos de} & = \lsg{weit weg von} \\[0.5em]
\es{a la izquierda} & = \lsg{nach links} & \es{al lado de} & = \lsg{neben} \\[0.5em]
\es{Sigue...} & = Geh weiter... & \es{enfrente de} & = \lsg{gegenüber von} \\[0.5em]
\end{tabular}

\end{merkkasten}

\vspace{15pt}

% ===============================================================
% AUFGABE 1: Orte zuordnen
% ===============================================================

\begin{aufgabe}{1 -- Ordne zu: Orte in der Stadt}
Verbinde die spanischen Orte mit der deutschen Bedeutung. \punktebox{4}
\end{aufgabe}

\begin{center}
\begin{tabular}{rcl}
\es{el parque} & \lsg{---} & das Kino \lsg{$\rightarrow$ el parque = der Park} \\[0.8em]
\es{la tienda} & \lsg{---} & der Park \lsg{$\rightarrow$ la tienda = das Geschäft} \\[0.8em]
\es{el cine} & \lsg{---} & die Apotheke \lsg{$\rightarrow$ el cine = das Kino} \\[0.8em]
\es{la farmacia} & \lsg{---} & das Geschäft \lsg{$\rightarrow$ la farmacia = die Apotheke} \\
\end{tabular}
\end{center}

\vspace{10pt}

% ===============================================================
% AUFGABE 2: Verben konjugieren
% ===============================================================

\begin{aufgabe}{2 -- Ergänze die Verben ir und estar}
Schreibe die richtige Form. \punktebox{4}
\end{aufgabe}

\noindent
a) Yo \lsg{voy} (ir) al supermercado. \hfill
b) Tú \lsg{estás} (estar) en el parque. \\[0.8em]
c) Ella \lsg{va} (ir) a la farmacia. \hfill
d) Nosotros \lsg{estamos} (estar) en la plaza. \\

\vspace{10pt}

% ===============================================================
% AUFGABE 3: Lückentext Wegbeschreibung
% ===============================================================

\begin{aufgabe}{3 -- Wegbeschreibung: Ergänze die Lücken}
Wähle aus: \es{todo recto}, \es{a la derecha}, \es{a la izquierda}, \es{cerca de}, \es{enfrente de}, \es{al lado de} \punktebox{6}
\end{aufgabe}

\noindent
a) -- ¿Dónde está el banco?\\
\hspace*{0.5cm} -- El banco está \lsg{al lado de} la iglesia. (neben)

\vspace{0.8em}

\noindent
b) -- ¿Cómo llego al cine?\\
\hspace*{0.5cm} -- Sigue \lsg{todo recto} y luego gira \lsg{a la derecha}. (geradeaus, nach rechts)

\vspace{0.8em}

\noindent
c) -- ¿Está lejos la farmacia?\\
\hspace*{0.5cm} -- No, está \lsg{cerca de} el hospital. (in der Nähe von)

\vspace{0.8em}

\noindent
d) -- ¿Dónde está el museo?\\
\hspace*{0.5cm} -- Está \lsg{enfrente de} el parque. (gegenüber von)

\vspace{10pt}

% ===============================================================
% AUFGABE 4: Dialog schreiben
% ===============================================================

\begin{aufgabe}{4 -- Schreibe einen Dialog: Wegbeschreibung}
A fragt nach dem Weg zum Kino. B beschreibt den Weg. \punktebox{6}
\end{aufgabe}

\noindent
\textbf{A:} \lsg{Perdona, ¿dónde está el cine? / ¿Cómo llego al cine?} (2P)

\vspace{0.8em}

\noindent
\textbf{B:} \lsg{Sigue todo recto y luego gira a la izquierda.} (2P)

\vspace{0.8em}

\noindent
\textbf{A:} ¿Está lejos?

\vspace{0.8em}

\noindent
\textbf{B:} \lsg{No, está cerca. / No, está al lado del parque.} (1P)

\vspace{0.8em}

\noindent
\textbf{A:} ¡Gracias!

\vspace{0.8em}

\noindent
\textbf{B:} \lsg{¡De nada!} (1P)

% ===============================================================
% PUNKTEÜBERSICHT
% ===============================================================

\vspace{20pt}
\hrule
\vspace{10pt}

\begin{flushright}
\textbf{Punkte:} \lsg{\maxpunkte} / \maxpunkte
\end{flushright}

\end{document}
